\section*{Chapter \ref{ch:b3:complex-spectra-magnetism} --- Electronic Spectra and Magnetism of Complexes}

\begin{solution}{exo:b3:complex-spectra-magnetism:1}
D: E$_g$ + T$_{2g}$ ($2 + 3 = 5$); F: A$_{2g}$ + T$_{1g}$ + T$_{2g}$ ($1 + 3 + 3 = 7$); G: A$_{1g}$ + E$_g$ + T$_{1g}$ + T$_{2g}$
($1 + 2 + 3 + 3 = 9$).
\end{solution}

\begin{solution}{exo:b3:complex-spectra-magnetism:2}
$d^1$ $^2$D $\to$ $^2$T$_{2g}$; $d^2$ $^3$F $\to$ $^3$T$_{1g}$; $d^3$ $^4$F $\to$ $^4$A$_{2g}$; $d^4$ $^5$D $\to$ $^5$E$_g$; $d^5$ $^6$S $\to$ $^6$A$_{1g}$; $d^6$ $^5$D
$\to$ $^5$T$_{2g}$; $d^7$ $^4$F $\to$ $^4$T$_{1g}$; $d^8$ $^3$F $\to$ $^3$A$_{2g}$; $d^9$ $^2$D $\to$ $^2$E$_g$.
\end{solution}

\begin{solution}{exo:b3:complex-spectra-magnetism:3}
High-spin \ce{Mn^{2+}} ($n = 5$) $5.92\,\mu_B$, \ce{Fe^{2+}} (4) $4.90\,\mu_B$, \ce{Co^{2+}} (3) $3.87\,\mu_B$; low-spin \ce{Fe^{2+}} and
\ce{Co^{3+}} ($t_{2g}^6$): 0, diamagnetic.
\end{solution}

\begin{solution}{exo:b3:complex-spectra-magnetism:4}
$\ce{[Mn(H2O)6]^{2+}}$ (spin- and Laporte-forbidden, almost colourless) $<$ $\ce{[Co(H2O)6]^{2+}}$
(Laporte-forbidden, vibronic) $<$ $\ce{[CoCl4]^{2-}}$ (no centre of symmetry) $<$ \ce{MnO4-} (allowed
charge transfer).
\end{solution}

\begin{solution}{exo:b3:complex-spectra-magnetism:5}
8500: $^3$A$_{2g}\to{}^3$T$_{2g}$, so $\Delta_{\mathrm o} = \qty{8500}{cm^{-1}}$; 14100: $^3$T$_{1g}$(F); 24900: $^3$T$_{1g}$(P).
$15B = 39\,000 - 3 \times 8500$, $B = \qty{900}{cm^{-1}}$, $\beta = 900/1056 = 0.85$.
\end{solution}

\begin{solution}{exo:b3:complex-spectra-magnetism:6}
$\Delta_{\mathrm o} = \qty{17400}{cm^{-1}}$; solving $7.5B + 1.5\Delta_{\mathrm o} - \frac12\sqrt{225B^2 - 18B\Delta_{\mathrm o} + \Delta_{\mathrm o}^2} = 24\,600$: $B = \qty{729}{cm^{-1}}$,
$\beta = 0.79$. $\nu_3 = \qty{38500}{cm^{-1}}$ (\qty{260}{nm}), in the ultraviolet.
\end{solution}

\begin{solution}{exo:b3:complex-spectra-magnetism:7}
The $d$ electrons are more delocalised onto the oxide ions of the gems than onto
water molecules: the Cr--O bonds there are more covalent.
\end{solution}

\begin{solution}{exo:b3:complex-spectra-magnetism:8}
Strong ($e_g$ unevenly filled): $d^9$, high-spin $d^4$, low-spin $d^7$. None: $d^3$ ($t_{2g}^3$, A term) and
low-spin $d^6$ ($t_{2g}^6$). Weak: high-spin $d^6$ ($t_{2g}^4e_g^2$, T term).
\end{solution}

\begin{solution}{exo:b3:complex-spectra-magnetism:9}
Octahedral high-spin $d^7$: ground term $^4$T$_{1g}$, a T term with an \omterm{def:b3:complex-spectra-magnetism:moment}{orbital contribution},
moment well above the spin-only $3.87\,\mu_B$. Tetrahedral $d^7$ ($\equiv$ octahedral $d^3$): ground
term $^4$A$_2$, orbital moment quenched, close to spin-only (a little above, by
mixing with excited T terms).
\end{solution}

\begin{solution}{exo:b3:complex-spectra-magnetism:10}
$\mu_{\mathrm{eff}} = 2.828\sqrt{1.87} = 3.87\,\mu_B$: three unpaired electrons ($S = \frac32$).
\end{solution}

\begin{solution}{exo:b3:complex-spectra-magnetism:11}
$\Delta\chi_v = 3 \times 25.0/\num{400e6} = \num{1.88e-7}$; $\chi_{\mathrm m} = \num{1.88e-7}/\qty{10.0}{mol.m^{-3}} = \qty{1.88e-8}{m^3.mol^{-1}}$, i.e.
\qty{1.49e-3}{cm^3.mol^{-1}}; $\chi_{\mathrm m}T = 0.445$, $\mu_{\mathrm{eff}} = 1.89\,\mu_B$: one unpaired electron.
\end{solution}

\begin{solution}{exo:b3:complex-spectra-magnetism:12}
$T_{1/2} = 15\,000/75 = \qty{200}{K}$. At \qty{250}{K}: $\Delta G = 15\,000 - 250 \times 75 = \qty{-3750}{J/mol}$, $K = \eu^{1.80} = 6.07$,
$x_{\mathrm{HS}} = 0.86$; $\chi_{\mathrm m}T = 0.86 \times 3.00 = \qty{2.6}{cm^3.K.mol^{-1}}$.
\end{solution}

\begin{solution}{pb:b3:complex-spectra-magnetism:1}
\textbf{1.} [Ar]$3d^3$.
\textbf{2.} $^4$F.
\textbf{3.} $^4$A$_{2g}$ + $^4$T$_{2g}$ + $^4$T$_{1g}$.
\textbf{4.} $^4$A$_{2g}$, the term of $t_{2g}^3$, all three electrons in the lower orbitals.
\textbf{5.} $^4$T$_{1g}$(P).
\textbf{6.} $^4$A$_{2g}\to{}^4$T$_{2g}$, $\to{}^4$T$_{1g}$(F), $\to{}^4$T$_{1g}$(P).
\textbf{7.} 561 and \qty{414}{nm}.
\textbf{8.} 597 and \qty{434}{nm}.
\textbf{9.} Ruby absorbs yellow-green and violet and transmits red (with a little
blue); emerald's bands, shifted to the red, absorb orange-red as well, and green
is transmitted.
\textbf{10.} Ruby \qty{17820}{cm^{-1}}, emerald \qty{16740}{cm^{-1}}.
\textbf{11.} $^4$T$_{2g}$ ($t_{2g}^2e_g$) and $^4$A$_{2g}$ ($t_{2g}^3$) have the same electron-repulsion energy; their
difference is the one-electron promotion energy $\Delta_{\mathrm o}$.
\textbf{12.} $B = (14\,305 - 547)/15 = \qty{917}{cm^{-1}}$.
\textbf{13.} \qty{614}{cm^{-1}}.
\textbf{14.} \qty{618}{cm^{-1}}.
\textbf{15.} 0.67 for both.
\textbf{16.} 29.0 and 27.1.
\textbf{17.} $\nu_3 = \qty{38500}{cm^{-1}}$ (\qty{260}{nm}) for ruby, in the ultraviolet, where charge-transfer
absorption hides it.
\textbf{18.} The octahedral site of beryl is larger: longer Cr--O bonds give a smaller
splitting, which varies steeply with the distance.
\textbf{19.} $10^7/695 = \qty{14390}{cm^{-1}}$, $23.4B$.
\textbf{20.} $^2$E$_g$ depends mainly on $C$: the higher measured level means $C/B$ is larger,
about 5.3 by the same diagonalisation.
\textbf{21.} $^2$E$_g$ and $^4$A$_{2g}$ both belong to $t_{2g}^3$, differing only by a spin flip: the bonding,
hence the geometry, is the same in both, and the transition has no vibrational
progression (a 0--0 line).
\textbf{22.} It is spin-forbidden.
\textbf{23.} Its energy depends on $B$ and $C$, almost not on $\Delta_{\mathrm o}$ (a flat line on the diagram),
and $B$ is nearly the same in the two gems.
\textbf{24.} Ions pumped into the broad bands relax and accumulate in the long-lived
$^2$E$_g$ level, so that more ions can be in it than in the ground level: a \omterm{def:b3:partition-functions:boltzmann}{population}
inversion.
\textbf{25.} \textbf{$\Delta_{\mathrm o}(\text{ruby}) - \Delta_{\mathrm o}(\text{emerald}) = \qty{1080}{cm^{-1}}$}: the whole difference between red and
green.
\end{solution}
